% ================================================================
%  sesion01.tex  —  Sesión 1 de 20
%  ¿Qué es una guía turística matemática?
%  Almería en Cifras y Formas · Matemáticas 1.º ESO
%  Compilar con:  pdflatex sesion01.tex  (dos pasadas)
% ================================================================
\documentclass[aspectratio=169, 10pt, spanish]{beamer}
\usepackage{almeria-beamer}

% ---------------------------------------------------------------
%  DATOS DE LA SESIÓN
% ---------------------------------------------------------------
\renewcommand{\SessionNumber}{01}
\renewcommand{\SessionTitle}{¿Qué es una guía turística matemática?}
\renewcommand{\SessionName}{Sesión 01}
\renewcommand{\SessionObjective}{%
  Identificar qué tipo de información puede tener carácter
  matemático en una ciudad y distinguir datos descriptivos
  de datos cuantitativos usando Almería como contexto.}
\renewcommand{\BlockName}{Bloque 1: Mapa y coordenadas}

% ================================================================
\begin{document}

% ---------------------------------------------------------------
%  PORTADA
% ---------------------------------------------------------------
\begin{frame}[plain]
  \titlepage
\end{frame}

% ---------------------------------------------------------------
%  ÍNDICE DE LA SESIÓN
% ---------------------------------------------------------------
\begin{frame}{Índice de la sesión}
  \begin{columns}[t]
    \begin{column}{0.5\textwidth}
      \begin{enumerate}
        \item ¿Qué es una guía turística?
        \item ¿Qué la hace \textit{matemática}?
        \item Tipos de información: descriptiva vs.\ cuantitativa
        \item Datos matemáticos ocultos en Almería
        \item Ejemplo paso a paso: el Cable Inglés
      \end{enumerate}
    \end{column}
    \begin{column}{0.5\textwidth}
      \begin{enumerate}
        \setcounter{enumi}{5}
        \item Ejercicios multinivel (Bloom)
        \item Evidencia del día
        \item Error frecuente
        \item Resumen
      \end{enumerate}
    \end{column}
  \end{columns}
\end{frame}

% ---------------------------------------------------------------
\seccionframe{1. ¿Qué es una guía turística?}
% ---------------------------------------------------------------

\begin{frame}{¿Qué es una guía turística?}
  \begin{columns}[c]
    \begin{column}{0.52\textwidth}
      \begin{definicionbox}{Guía turística}
        Publicación (papel o digital) que describe los
        lugares de interés de una ciudad o región:
        monumentos, museos, rutas, horarios, precios\ldots
      \end{definicionbox}
      
      \begin{itemize}
        \item Cuenta \textbf{qué} hay y \textbf{cómo llegar}.
        \item Usa palabras, fotos y mapas.
        \item Va dirigida al \textbf{turista}.
      \end{itemize}
      
      \begin{notabox}[Ejemplos conocidos]
        
        Lonely Planet · Guía Repsol · Tripadvisor
      \end{notabox}
    \end{column}
    \begin{column}{0.44\textwidth}
      \begin{center}
        \begin{tikzpicture}[scale=0.9]
          % Portada de libro estilizada
          \fill[AlmAzul, rounded corners=4pt]
            (0,0) rectangle (3.2,4.4);
          \fill[AlmAmbar]
            (0.15,3.6) rectangle (3.05,4.2);
          \node[text=AlmBlanco, font=\bfseries] at (1.6,3.9)
            {ALMERÍA};
          \fill[AlmAzulC!40]
            (0.15,2.1) rectangle (3.05,3.5);
          \node[text=AlmAzul, font=] at (1.6,2.8)
            {mapa de la ciudad};
          \foreach \y in {1.2,0.85,0.5}{
            \fill[AlmBlanco!30, rounded corners=1pt]
              (0.3,\y) rectangle (2.8,\y+0.22);
          }
          \node[text=AlmAmbar, font=\bfseries] at (1.6,0.05)
            {Guía Turística};
        \end{tikzpicture}
      \end{center}
    \end{column}
  \end{columns}
\end{frame}

% ---------------------------------------------------------------
\seccionframe{2. ¿Qué la hace \textit{matemática}?}
% ---------------------------------------------------------------

\begin{frame}{¿Qué hace matemática a una guía?}
  \begin{columns}[c]
    \begin{column}{0.55\textwidth}
      \begin{teoriabox}
        Una guía turística \textbf{matemática} incorpora,
        además de texto descriptivo, datos numéricos,
        medidas y representaciones cuantitativas que
        permiten \textit{analizar} y \textit{comparar}
        la realidad de la ciudad.
      \end{teoriabox}
      {8pt}
      ¿Qué puede aparecer?
      \begin{itemize}
        \item \textbf{Coordenadas} de monumentos
        \item \textbf{Distancias} entre puntos turísticos
        \item \textbf{Alturas}, superficies, aforos
        \item \textbf{Gráficas} de visitantes o temperatura
        \item \textbf{Estadísticas} históricas
      \end{itemize}
    \end{column}
    \begin{column}{0.42\textwidth}
      \begin{tikzpicture}[every node/.style={font=}]
        % Diagrama Venn
        \fill[AlmAzulC!25, opacity=0.7] (-0.6,0) ellipse (2cm and 1.4cm);
        \fill[AlmAmbar!35, opacity=0.7] (0.6,0) ellipse (2cm and 1.4cm);
        \node[text=AlmAzul, align=center] at (-1.4,0) {Guía\\clásica};
        \node[text=AlmAzul] at (2.0,0) {Matemáticas};
        \node[text=AlmNegro, align=center] at (0,0)
          {Guía\\matemática};
        \node[text=AlmGris, font=, align=center] at (-1.5,-0.55)
          {fotos\\textos};
        \node[text=AlmGris, font=, align=center] at (1.6,-0.55)
          {cálculos\\gráficas};
        \node[text=AlmAzul, font=, align=center] at (0,0.55)
          {datos\\medidas};
      \end{tikzpicture}
    \end{column}
  \end{columns}
\end{frame}

% ---------------------------------------------------------------
\seccionframe{3. Información descriptiva vs.\ cuantitativa}
% ---------------------------------------------------------------

\begin{frame}{Dos tipos de información}
  \begin{columns}[T]
    \begin{column}{0.48\textwidth}
      \begin{definicionbox}{Información descriptiva}
        \begin{itemize}
          \item Usa palabras, no números.
          \item Describe \textbf{cómo es} algo.
          \item Ejemplo: «La Alcazaba es una fortaleza árabe
                de impresionante belleza.»
        \end{itemize}
      \end{definicionbox}
    \end{column}
    \begin{column}{0.48\textwidth}
      \begin{definicionbox}{Información cuantitativa}
        \begin{itemize}
          \item Usa \textbf{números y unidades}.
          \item Dice \textbf{cuánto} o \textbf{cuántos}.
          \item Ejemplo: «La Alcazaba ocupa
                $43\,000\ \text{m}^2$ y
                recibe $300\,000$ visitantes/año.»
        \end{itemize}
      \end{definicionbox}
    \end{column}
  \end{columns}
  
  \begin{center}
    \begin{tabular}{lll}
      \toprule
      \textbf{Monumento} & \textbf{Descriptivo} & \textbf{Cuantitativo}\\
      \midrule
      Alcazaba   & Fortaleza islámica medieval & $43\,000\ \text{m}^2$, s.\ X\\
      Cable Inglés & Viaducto de hierro único   & $1{,}1\ \text{km}$, 1904\\
      Catedral   & Estilo gótico-renacentista   & 55 m altura, 1524\\
      \bottomrule
    \end{tabular}
  \end{center}
\end{frame}

% ---------------------------------------------------------------
\seccionframe{4. Datos matemáticos ocultos en Almería}
% ---------------------------------------------------------------

\begin{frame}{Datos matemáticos ocultos}
  \begin{columns}[T]
    \begin{column}{0.6\textwidth}
      Cualquier elemento urbano esconde información
      matemática. Hay que saber \textbf{hacer preguntas}:
      
      \begin{tabular}{@{}p{2.8cm}p{5cm}@{}}
        \toprule
        \textbf{Lugar} & \textbf{Preguntas matemáticas}\\
        \midrule
        Alcazaba
          & ¿Cuántos metros de muralla tiene?\\[2pt]
        Puerto de Almería
          & ¿Qué área ocupa? ¿Cuántos barcos atraca?\\[2pt]
        Paseo Marítimo
          & ¿Cuántos km mide de extremo a extremo?\\[2pt]
        Catedral
          & ¿Cuánto tarda en dar vuelta la sombra de la torre?\\
        \bottomrule
      \end{tabular}
    \end{column}
    \begin{column}{0.37\textwidth}
      \begin{notabox}[¿Sabías que\ldots?]
        
        El \textbf{Puerto de Almería} ocupa
        \textbf{35 ha}, lo que equivale a
        $35\times 10\,000 = 350\,000\ \text{m}^2$.
        ¡Más de 49 campos de fútbol!
      \end{notabox}
      
      \begin{tikzpicture}[scale=0.75]
        \fill[AlmAzulC!30, rounded corners=3pt]
          (0,0) rectangle (3.8,2.2);
        \fill[AlmVerde!40]
          (0.2,0.2) rectangle (1.6,0.6);
        \node[font=, color=AlmAzul] at (0.9,0.4)
          {campo fútbol};
        \node[font=\bfseries, color=AlmAzul] at (1.9,1.2)
          {$35\ \text{ha}$};
        \node[font=, color=AlmGris] at (1.9,0.7)
          {Puerto};
      \end{tikzpicture}
    \end{column}
  \end{columns}
\end{frame}

% ---------------------------------------------------------------
\seccionframe{5. Ejemplo: el Cable Inglés}
% ---------------------------------------------------------------

\begin{frame}{El Cable Inglés — Ejemplo paso a paso}

  \framesubtitle{Longitud 1,1 km · Velocidad media de paseo 5 km/h}
  
        \begin{ejemplobox}
          
          El Cable Inglés es un viaducto de hierro de
          $1{,}1\ \text{km}$ de longitud que cruza la
          desembocadura de la Rambla.\\[4pt]
          \textbf{Pregunta:} ¿Cuánto tarda en cruzarlo
          un peatón a $5\ \text{km/h}$?
        \end{ejemplobox}
  \begin{columns}[c]
    \begin{column}{0.5\textwidth}
      
      \begin{pasobox}[1]
        
        \textbf{Datos:}
        $d = 1{,}1\ \text{km}$;\quad
        $v = 5\ \text{km/h}$
      \end{pasobox}
      
            \begin{pasobox}[2]
              
              \textbf{Fórmula tiempo--distancia--velocidad:}
              \[
                t = \frac{d}{v}
              \]
            \end{pasobox}
    \end{column}
    \begin{column}{0.46\textwidth}
      
      \begin{pasobox}[3]
        
        \textbf{Sustitución:}
        \[
          t = \frac{1{,}1\ \text{km}}{5\ \text{km/h}}
            = 0{,}22\ \text{h}
        \]
      \end{pasobox}
      
      \begin{pasobox}[4]
        
        \textbf{Conversión a minutos:}
        \[
          t = 0{,}22 \times 60 \approx \mathbf{13\ \text{min}}
        \]
      \end{pasobox}
    \end{column}
  \end{columns}
\end{frame}

% ---------------------------------------------------------------
\seccionframe{6. Ejercicios multinivel — Bloom}
% ---------------------------------------------------------------

\begin{frame}{Ejercicios multinivel}
  \begin{columns}[T]
    \begin{column}{0.31\textwidth}
      \begin{recordarbox}
        
        \textbf{Define} con tus palabras qué es
        una guía turística.\par
        Escribe \textbf{3 monumentos} de Almería
        y anota 2 datos matemáticos que podrías
        buscar de cada uno.
      \end{recordarbox}
    \end{column}
    \begin{column}{0.31\textwidth}
      \begin{comprenderbox}
        
        El Paseo de Almería mide $370\ \text{m}$
        y se recorre a $4\ \text{km/h}$.\par
        \textbf{Explica} qué tipo de información es
        cada dato y calcula cuántos minutos dura
        el paseo.
      \end{comprenderbox}
    \end{column}
    \begin{column}{0.31\textwidth}
      \begin{analizarbox}
        
        Observa tres monumentos de tu elección.\par
        \textbf{Extrae} al menos 2 preguntas
        matemáticas que podrían incluirse en
        una guía turística para cada uno.\par
        ¿Qué tipo de matemáticas necesitarías
        para responderlas?
      \end{analizarbox}
    \end{column}
  \end{columns}
\end{frame}

% ---------------------------------------------------------------
%  EVIDENCIA
% ---------------------------------------------------------------
\begin{frame}{Evidencia del día}
  \begin{evidenciabox}
    \begin{itemize}
      \item Entrega tu \textbf{tarjeta de trabajo} con
            al menos \textbf{5 elementos} que incluirás
            en vuestra guía turística matemática.
      \item Para cada elemento indica:
            \begin{enumerate}
              \item Nombre del monumento o lugar
              \item Un dato \textit{descriptivo} y uno \textit{cuantitativo}
              \item Una pregunta matemática que se podría resolver
            \end{enumerate}
      \item Incluye las operaciones intermedias, no solo el resultado.
    \end{itemize}
  \end{evidenciabox}
\end{frame}

% ---------------------------------------------------------------
%  ERROR FRECUENTE
% ---------------------------------------------------------------
\begin{frame}{¡Cuidado con este error!}
  \begin{errorbox}
    % --- CORRECCIÓN: Uso de minipage en lugar de columns ---
    \begin{minipage}[t]{0.48\textwidth}
      \correcto\ \textbf{Correcto} \\[4pt]
      Las matemáticas sirven para \textbf{describir,
      medir y comparar} la realidad:\par
      «La Alcazaba recibe
      $\mathbf{300\,000}$ visitas/año
      y tiene $\mathbf{43\,000\ \text{m}^2}$.»
    \end{minipage}%
    \hfill%
    \begin{minipage}[t]{0.48\textwidth}
      \incorrecto\ \textbf{Error frecuente}\\[4pt]
      Pensar que las matemáticas son
      \textbf{solo cálculo}:\par
      «Las matemáticas no tienen nada que ver
      con describir ciudades.»
    \end{minipage}
    % -------------------------------------------------------
  \end{errorbox}
  \begin{notabox}[Recuerda]
    
    Todo lo que se puede \textbf{medir, contar, localizar
    o comparar} en una ciudad es material matemático.
  \end{notabox}
\end{frame}

% ---------------------------------------------------------------
%  RESUMEN
% ---------------------------------------------------------------
\begin{frame}{Resumen de la sesión 01}
  \begin{resumenbox}
    \begin{itemize}
      \item Una \textbf{guía turística matemática} combina
            descripciones con datos numéricos, medidas y
            representaciones cuantitativas.
      \item Los datos pueden ser \textbf{descriptivos}
            (palabras) o \textbf{cuantitativos} (números).
      \item Cualquier monumento o espacio urbano esconde
            información matemática si sabemos \textbf{hacer
            las preguntas correctas}.
      \item Hemos aplicado la fórmula
            $t = \dfrac{d}{v}$ al Cable Inglés.
    \end{itemize}
  \end{resumenbox}
  
  \begin{center}
    {\color{AlmAzulM}
    \faArrowCircleRight\enskip
    \textbf{Próxima sesión:}
    ¿Qué se puede medir y representar en una ciudad?}
  \end{center}
\end{frame}

\end{document}


